\documentclass{report}
\usepackage{amsfonts, amsmath}

\newcommand\N{{\mathbb N}}
\newcommand\Q{{\mathbb Q}}
\newcommand\R{{\mathbb R}}
\newcommand\Z{{\mathbb Z}}

\pagestyle{empty}
\begin{document}
\large
\centerline{\Huge \bf Analisi Matematica I}
\vskip.2cm
\centerline{\Huge  Terzo Appello (18-06-2002)}
\renewcommand{\thefootnote}{ } 
\footnote{\sl Avete 3 ore di tempo. Ogni esercizio
vale sei punti. La sufficienza si ottiene con un punteggio $\geq$
18. {\bf Solo le risposte chiaramente giustificate saranno prese in
considerazione.} {\scshape Elaborati illeggibili o confusi verranno ignorati.}
{\sffamily La
sintesi sar\`a particolarmente apprezzata.}}
\renewcommand{\thefootnote}{\arabic{footnote}} 

\vskip1cm
\begin{enumerate}
\item Dire quale dei seguenti numeri \`e il maggiore
\[
2^{10!}\;;\quad 1000000^{200}\;;\quad (1000!)^2.
\]
\item Dire per quali $q\in\N$ si ha che
\[
\frac{q+1}{q-1}\in\Z.
\]
\item Mostrare che, per ogni $q\in\N$, si ha che
\[
\sum_{n=q+1}^\infty\frac 1{n!}\leq \frac1{q!q},
\]
e che, se $q>3$, $p\in\N$,
\[
\left|\frac pq-\sum_{n=0}^q\frac 1{n!}\right|\geq \frac 1{q!}.
\]
Si usino queste due disuguaglianze per dimostrare che $e\not\in\Q$.

\item Si studi la convergenza della serie e si calcoli il limite
\[
\sum_{n=1}^\infty \left(1-\frac 1n\right)^{n^2}\;;\quad
\lim_{n\to\infty} n\left\{\sqrt{n^2+1}-n\right\}.
\]
\item \label{ex:5}Si determini il dominio e si disegni il grafico della funzione
\[
f(x)=\sqrt{2-\sqrt{4-x^2}}.
\]
\item Sia $\ell_0=2$ e si definisca, per ogni $n\geq 0$,
$\ell_{n+1}=f(\ell_n)$ (dove $f$ \`e la funzione definita nella
domanda \eqref{ex:5}). Si dimostri che la successione $a_n:=2^n\ell_n$
\`e limitata.
\end{enumerate}
\end{document}
